%% kapa: exam=1 round=28 year=2017 period=2 form=A subject=accounting subject_ko=회계학 seq=31 number=71
%% source: Q-Net 감정평가사 기출문제 (cst00309) artlSeq=5199603 file=extracted/1008/5199603/2017년 제28회 감정평가사 1차시험 문제지/2교시 A형.pdf
%% conversion: kordoc 4.24.51 -> tools/gen1.py (auto); review=proofread-sonnet; flags=none
\begin{question}{71}
(주)대한은 제1공정에서 주산물 A, B와 부산물 C를 생산한다. 주산물 A와 부산물 C는 즉시 판매될 수 있으나, 주산물 B는 제2공정에서 추가가공을 거쳐 판매된다. 20x1년에 제1공정과 제2공정에서 발생된 제조원가는 각각 ￦150,000과 ￦60,000이었고, 제품별 최종 판매가치 및 판매비는 다음과 같다.
\par{\small\centering\begin{tabular}{|c|c|c|}
\hline
구분 & 최종 판매가치 & 판매비 \\ \hline
A & ￦100,000 & ￦2,000 \\ \hline
B & 180,000 & 3,000 \\ \hline
C & 2,000 & 600 \\ \hline
\end{tabular}\par}
(주)대한은 주산물의 매출총이익률이 모두 동일하게 되도록 제조원가를 배부하며, 부산물은 판매시점에 최초로 인식한다. 주산물 A의 총제조원가는? (단, 기초 및 기말 재고자산은 없다.) \par
\begin{choices}
  \item ￦74,500
  \item ￦75,000
  \item ￦76,000
  \item ￦77,500
  \item ￦78,000
\end{choices}
\end{question}
\answer{2}
